% GENERATED FILE. Edit canonical lessons, then run npm run generate:books.
% canonical-source-hash: fnv1a64:2bdd2309814a986a
% canonical-lessons: SA-C44-white, SA-C44-dark, SA-C44-yellow, SA-C44-green, SA-C44-blue

\chapter{Colours}
\label{ch:sa-colours}
\hlchaptermodality{44}

\begin{chapteropening}
\textbf{By the end of this chapter:} \emph{I can say white, dark, yellow, green and blue in Sanskrit.}

Complete SA-C44-blue and demonstrate the chapter promise.
\end{chapteropening}

\section[śvetaḥ]{\sk{श्वेतः} (śvetaḥ) — white}

\label{lesson:SA-C44-white}



\begin{quote}
Before the new word: what did \emph{saralaḥ} mean?
\end{quote}

\subsection*{You'll want to know: \sk{श्वेतः}}
\textbf{\sk{श्वेतः}} (\emph{śvetaḥ}) — \textquotedblleft{}white\textquotedblright{}.

\textbf{\sk{श्वेतः}} is white, and bright along with it.

English \emph{white} is the same word by the Germanic road, and Russian \emph{svet}, \textquotedblleft{}light\textquotedblright{}, is another of its brothers — brightness and whiteness were one idea for a long time.

\sk{हिमम्} is \textbf{\sk{श्वेतम्}}, and so is \sk{दधि}.

\begin{grammarlens}[title={What you've built}]
White, and the brightness that comes with it.
\end{grammarlens}

\subsection*{Your turn}
\begin{itemize}
\raggedright
  \item \emph{Say it:} \emph{śvetaḥ}
  \item \emph{Say it:} it once more, slowly
  \item \emph{Say it:} \emph{saralaḥ}, then \emph{śvetaḥ}
  \item \emph{Recall:} say \emph{kadācit}, then read \textbf{\sk{आतिथ्यम्}}
\end{itemize}

\begin{tcolorbox}[breakable,colback=teal!4,colframe=teal!35!black,title={Before you move on}]
What does \sk{श्वेतः} mean? (\textquotedblleft{}white\textquotedblright{}.) Say it once more.
\end{tcolorbox}
\section[śyāmaḥ]{\sk{श्यामः} (śyāmaḥ) — dark, black}

\label{lesson:SA-C44-dark}



\begin{quote}
Before the new word: what did \emph{śvetaḥ} mean?
\end{quote}

\subsection*{You'll want to know: \sk{श्यामः}}
\textbf{\sk{श्यामः}} (\emph{śyāmaḥ}) — \textquotedblleft{}dark, black\textquotedblright{}.

\textbf{\sk{श्यामः}} is dark — black, but also the blue-black of a rain cloud and the deep green of standing water.

It is the darkness of a thundercloud rather than of soot, which is why the god whose name means \textquotedblleft{}the dark one\textquotedblright{} is also called \emph{Śyām} with affection.

No European cousin has been agreed. The word stayed home, and \emph{Śyām} and \emph{Śyāmā} are ordinary names in India today.

\begin{grammarlens}[title={What you've built}]
Dark, with a rain cloud in it rather than ash.
\end{grammarlens}

\subsection*{Your turn}
\begin{itemize}
\raggedright
  \item \emph{Say it:} \emph{śyāmaḥ}
  \item \emph{Say it:} it once more, slowly
  \item \emph{Say it:} \emph{śvetaḥ}, then \emph{śyāmaḥ}
  \item \emph{Recall:} read \textbf{\sk{सहसा}}, then say \emph{vinayaḥ}
\end{itemize}

\begin{tcolorbox}[breakable,colback=teal!4,colframe=teal!35!black,title={Before you move on}]
What does \sk{श्यामः} mean? (\textquotedblleft{}dark, black\textquotedblright{}.) Say it once more.
\end{tcolorbox}
\section[pītaḥ]{\sk{पीतः} (pītaḥ) — yellow}

\label{lesson:SA-C44-yellow}



\begin{quote}
Before the new word: what did \emph{śyāmaḥ} mean?
\end{quote}

\subsection*{You'll want to know: \sk{पीतः}}
\textbf{\sk{पीतः}} (\emph{pītaḥ}) — \textquotedblleft{}yellow\textquotedblright{}.

\textbf{\sk{पीतः}} is yellow — turmeric, ripe grain, a marigold, the silk a bridegroom wears.

Its older force is \textquotedblleft{}made yellow\textquotedblright{}, so a \textbf{\sk{पीतम्}} \sk{वस्त्रम्} is a cloth that has been dyed rather than one that grew that colour.

Hindi \emph{pīlā} is the same word with an ending on it, and it is the yellow of every wall and every taxi.

\begin{grammarlens}[title={What you've built}]
Yellow, understood as a thing that was dyed.
\end{grammarlens}

\subsection*{Your turn}
\begin{itemize}
\raggedright
  \item \emph{Say it:} \emph{pītaḥ}
  \item \emph{Say it:} it once more, slowly
  \item \emph{Say it:} \emph{śyāmaḥ}, then \emph{pītaḥ}
  \item \emph{Recall:} say \emph{prāyaḥ}, then read \textbf{\sk{नम्रता}}
\end{itemize}

\begin{tcolorbox}[breakable,colback=teal!4,colframe=teal!35!black,title={Before you move on}]
What does \sk{पीतः} mean? (\textquotedblleft{}yellow\textquotedblright{}.) Say it once more.
\end{tcolorbox}
\section[haritaḥ]{\sk{हरितः} (haritaḥ) — green}

\label{lesson:SA-C44-green}



\begin{quote}
Before the new word: what did \emph{pītaḥ} mean?
\end{quote}

\subsection*{You'll want to know: \sk{हरितः}}
\textbf{\sk{हरितः}} (\emph{haritaḥ}) — \textquotedblleft{}green\textquotedblright{}.

\textbf{\sk{हरितः}} is green — and before it settled it covered yellow-green and gold as well.

English \emph{yellow} and \emph{gold} are its cousins, from \textbf{*ǵhel-}, a root that runs across the whole band from green to yellow. One family, one word, two colours pulled apart by different daughters.

\sk{तृणम्} is \textbf{\sk{हरितम्}}, and so is a young \sk{पत्रम्} on a \sk{तरुः}.

\begin{grammarlens}[title={What you've built}]
Green, from a root that could not decide between green and gold.
\end{grammarlens}

\subsection*{Your turn}
\begin{itemize}
\raggedright
  \item \emph{Say it:} \emph{haritaḥ}
  \item \emph{Say it:} it once more, slowly
  \item \emph{Say it:} \emph{pītaḥ}, then \emph{haritaḥ}
  \item \emph{Recall:} read \textbf{\sk{किञ्चित्}}, then say \emph{viśvāsaḥ}
\end{itemize}

\begin{tcolorbox}[breakable,colback=teal!4,colframe=teal!35!black,title={Before you move on}]
What does \sk{हरितः} mean? (\textquotedblleft{}green\textquotedblright{}.) Say it once more.
\end{tcolorbox}
\section[nīlaḥ]{\sk{नीलः} (nīlaḥ) — blue}

\label{lesson:SA-C44-blue}



\begin{quote}
Before the last word of the chapter: what did \emph{pītaḥ} mean, and what did \emph{haritaḥ} mean?
\end{quote}

\subsection*{You'll want to know: \sk{नीलः}}
\textbf{\sk{नीलः}} (\emph{nīlaḥ}) — \textquotedblleft{}blue\textquotedblright{}.

\textbf{\sk{नीलः}} is blue, and dark blue at that — the blue of the \sk{समुद्रः} rather than of the \sk{आकाशः} at noon.

The indigo plant is \emph{nīlī}, and when Europe began buying that dye the word travelled with the cargo: Arabic \emph{nīl}, Portuguese \emph{anil}, and the \emph{aniline} English still uses for the dyes that replaced it.

So a Sanskrit colour word is sitting inside a modern chemistry term, carried there by trade.

\begin{grammarlens}[title={What you've built}]
Five colours: white, dark, yellow, green and blue.
\end{grammarlens}

\subsection*{Your turn}
\begin{itemize}
\raggedright
  \item \emph{Say it:} \emph{nīlaḥ}
  \item \emph{Say it:} it once more, slowly
  \item \emph{Say it:} all five in order, then \emph{śvetaḥ} and \emph{nīlaḥ} together
  \item \emph{Recall:} say \emph{bahu}, then read \textbf{\sk{सरलः}}
\end{itemize}

\begin{tcolorbox}[breakable,colback=teal!4,colframe=teal!35!black,title={Before you move on}]
What does \sk{नीलः} mean? (\textquotedblleft{}blue\textquotedblright{}.) Say it once more.
\end{tcolorbox}
